% Source pp. 10-19 (Lectures1-6.pdf).
\section{Quantization}

Before going to quantum theory, let us recall the Hamiltonian formalism of
classical mechanics and see how it looks for fields.

\subsection*{Hamiltonian Formalism}

In Lagrangian dynamics with the generalized coordinates $q^i$ one passes to
the Hamiltonian formalism through introducing the conjugate momenta
\[
p_i=\frac{\partial L(q,\dot q)}{\partial\dot q^i}\,,
\]
and then excluding $\dot q$ in favor of $p$, by solving these equations. The
Hamiltonian function $H$ appears as the Legendre transform
\[
H(p,q)=\sum_i p_i\,\dot q^i-L\,.
\]
The equations of motion then take the canonical form
\[
\dot p_i=\{p_i,H\}=-\frac{\partial H}{\partial q^i}\,,\qquad
\dot q^i=\{q^i,H\}=\frac{\partial H}{\partial p_i}\,,
\]
in terms of the Poisson brackets
\[
\{f,g\}:=\sum_i\frac{\partial f}{\partial q^i}\,\frac{\partial g}{\partial p_i}-\frac{\partial f}{\partial p_i}\,\frac{\partial g}{\partial q^i}\,.
\]
In a field theory we have continuously many generalized coordinates
$\varphi(\xx)$ (I assume here a scalar field, for simplicity), labeled by the
spatial coordinates $\xx$. In this case the conjugated momenta are defined as
\[
\pi(\xx)=\frac{\delta L}{\delta\dot\varphi(\xx)}\,,
\]
with the basic Poisson brackets
\[
\{\varphi(\xx),\pi(\xx')\}=\delta(\xx-\xx')\,.
\]
In the Klein-Gordon theory
\[
\pi(\xx)=\dot\varphi(\xx)\,,
\]
and the Hamiltonian is
\begin{equation}\label{eq:HKG}
H=\int d^3\xx\,\left[\frac12\,\pi^2(\xx)+\frac12\,\bigl(\nabla\varphi\bigr)^2(\xx)+\frac{m^2}{2}\,\varphi^2(\xx)\right].
\end{equation}
Of course, it coincides with the above expression for the energy,
$E=\int T^{00}\,d^3\xx$, expressed through $\pi$ rather than $\dot\varphi$.

\subsection*{Quantum Klein-Gordon field}

We can try now to develop quantum theory of the Klein-Gordon field, following
the usual rules of canonical quantization. One starts with the correspondence
\begin{center}
\begin{tabular}{ll}
Classical Theory & Quantum Theory\\[2pt]
Phase space & Hilbert space $\Ham$\\
$p_i,\,q^i$ & Hermitian operators $\hat p_i,\hat q^i:\ \Ham\to\Ham$\\
$H(p,q)$ & $\hat H(\hat p,\hat q):\ \Ham\to\Ham$
\end{tabular}
\end{center}
where the Hermitian operators $\hat q,\hat p$ must obey the canonical
commutation relations
\[
[\hat q^i,\hat p_j]=i\,\delta^i_j\,.
\]

\begin{remark}
In QM with finitely many degrees of freedom, in the Schroedinger's
quantization one can take the space of square-integrable functions $\Psi(q)$
for $\Ham$, and define the corresponding representation
\[
\hat q^i\,\Psi(q)=q^i\,\Psi(q)\,,\qquad
\hat p_i\,\Psi(q)=-i\,\frac{\partial}{\partial q^i}\,\Psi(q)\,,
\]
and $\hat H=H(\hat p,\hat q)$. The time evolution of the state $\Psi(q)$ is
governed by the Schroedinger equation
\[
i\,\frac{\partial}{\partial t}\,\Psi(q)=\hat H\,\Psi(q)\,.
\]
\end{remark}

In principle, one can follow this prescription and introduce the operators
$\hat\varphi(\xx)$, $\hat\pi(\xx)$ and demand that they satisfy the canonical
commutators in the form
\[
[\hat\varphi(\xx),\hat\pi(\xx')]=i\,\delta(\xx-\xx')\,.
\]
Furthermore, one can take the space of functionals $\Psi[\varphi(\xx)]$ as
the Hilbert space $\Ham$; it provides a formal representation of the above
canonical commutation relations through the identifications
\[
\hat\varphi(\xx)\,\Psi[\varphi]=\varphi(\xx)\,\Psi[\varphi]\,,\qquad
\hat\pi(\xx)\,\Psi[\varphi]=-i\,\frac{\delta}{\delta\varphi(\xx)}\,\Psi[\varphi]\,.
\]
The Hamiltonian operator then is
\[
\hat H=\int d^3\xx\,\left[-\frac12\,\frac{\delta^2}{\delta\varphi(\xx)^2}+\frac12\,\bigl(\nabla\varphi\bigr)^2+\frac{m^2}{2}\,\varphi^2\right],
\]
and it remains to find all solutions of the stationary Schroedinger equation
\[
\hat H\,\Psi[\varphi]=E\,\Psi[\varphi]\,.
\]

\begin{exercise}
Show that this Schroedinger equation admits formal solution of the form
\[
\Psi_0[\varphi]=\Const\,\exp\Bigl\{\int\varphi(\xx)\,G(\xx-\xx')\,\varphi(\xx')\,d^3\xx\,d^3\xx'\Bigr\}\,.
\]
Find the kernel $G(\xx-\xx')$. (At this point, do not worry when you discover
divergent integrals). Try to look for further solutions in the form
\[
\Psi_n[\varphi]\sim\int F(\xx_1,\dots,\xx_n)\,\varphi(\xx_1)\cdots\varphi(\xx_n)\,
\exp\Bigl\{\int\varphi(\xx)\,G(\xx-\xx')\,\varphi(\xx')\,d^3\xx\,d^3\xx'\Bigr\}\,.
\]
\end{exercise}

Although this straightforward approach, with certain refinements, would work
in the Klein-Gordon theory, it has a number of unpleasant features. One is
that the scalar product in $\Ham$ has to be defined through the functional
integral
\[
(\Psi_1,\Psi_2)\sim\int\Psi_1^*[\varphi]\,\Psi_2[\varphi]\,D[\varphi]\,,
\]
and certain care must be taken to give it a satisfactory definition (we are
going to use functional integrals, in somewhat different context, though).
This and other features of the above representation in terms of the
variational derivatives makes this approach somewhat cumbersome. Much shorter
route is based on explicit separation of variables, as we discuss below.

One observes that the Hamiltonian functional $H[\varphi]$ of the
Klein-Gordon theory is quadratic in the field variables $\varphi(\xx)$ and
their conjugate momenta $\pi(\xx)$; this means that the theory can be
understood as the \emph{collection of harmonic oscillators} (this observation
actually suggests the Gaussian form of $\Psi_0[\varphi]$ in the above
Exercise). It is then straightforward to make linear transformation of the
canonical variables $\varphi(\xx),\pi(\xx)$ which separates the variables into
a system of non-interacting harmonic oscillators, and then quantize the
individual oscillators. This transformation is described in detail in PS,
and I refer to \textbf{Sect.~2.3} there for the detailed calculation.

Instead, here I will follow equivalent routine based on special symmetries of
the KG theory. Recall that the Lagrangian density is
\[
\Lag=-\frac12\,\partial_\mu\varphi\,\partial^\mu\varphi-\frac{m^2}{2}\,\varphi^2\,,
\]
and observe that the transformation of the field
\[
\varphi(x)\to\varphi'(x)=\varphi(x)+f(x)\,,
\]
where $f(x)$ is an arbitrary solution of the KG equation
\[
\partial_\mu\partial^\mu f(x)-m^2\,f(x)=0\,,
\]
changes the Lagrangian density by a total derivative term
\[
\Lag(\varphi',\partial\varphi')=\Lag(\varphi,\partial\varphi)+\partial_\mu\tilde J^\mu_f\,,
\qquad\text{where}\qquad
\tilde J^\mu_f(x)=-\varphi(x)\,\partial^\mu f(x)
\]
(in writing this equation I have assumed that $f$ is infinitesimal and
neglected the terms $\sim f^2$). The situation meets the conditions of the
Noether's theorem. It follows that every solution $f$ of the KG equation
generates a conserved current $J^\mu_f$, such that
\[
\partial_\mu J^\mu_f=0\qquad\text{on shell.}
\]
Explicitly
\[
J^\mu_f=\frac{\partial\Lag}{\partial(\partial_\mu\varphi)}\,f-\tilde J^\mu_f=\varphi\,\partial^\mu f-\partial^\mu\varphi\,f\,.
\]
The corresponding Integrals of Motion are (at a time $t$)
\[
A_f=\int d^3\xx\,\Bigl(\dot\varphi(\xx,t)\,f(\xx,t)-\varphi(\xx,t)\,\dot f(\xx,t)\Bigr)
=\int d^3\xx\,\Bigl(\pi(\xx)\,f(\xx)-\varphi(\xx)\,\dot f(\xx)\Bigr)\,,
\]
where $\pi(\xx)=\dot\varphi(\xx)$ is the canonical momentum conjugate to
$\varphi(\xx)$. The integrals $A_f$ are Lorentz scalars. This becomes explicit
if one rewrites their definition in covariant form
\[
A_f=\int_\Sigma d\Sigma_\mu\,\bigl(\varphi\,\partial^\mu f-\partial^\mu\varphi\,f\bigr)\,,
\]
where $\Sigma$ is any space-like 3-surface (figure~\ref{fig:spacelike});
because of the continuity equation the currents $J^\mu_f$ obey, the above
integrals do not depend on specific choice of $\Sigma$.
\begin{figure}[ht]
\centering\figfile{spacelike}
\caption{A space-like 3-surface $\Sigma$. The integral $A_f$ over $\Sigma$
does not depend on the choice of $\Sigma$.}
\label{fig:spacelike}
\end{figure}

In classical theory one straightforwardly evaluates the Poisson brackets
\[
\{A_f,A_g\}=\int d^3\xx\,\Bigl(f(\xx,t)\,\dot g(\xx,t)-\dot f(\xx,t)\,g(\xx,t)\Bigr)
\]
(this expression also can be written in the covariant form, in terms of the
integral over arbitrary $\Sigma$). In quantum theory we should replace, as
usual,
\[
\{A,B\}\to-i\,[\hat A,\hat B]\,,
\]
so that
\[
[\hat A_f,\hat A_g]=-i\int d^3\xx\,\Bigl(\dot f(\xx)\,g(\xx)-f(\xx)\,\dot g(\xx)\Bigr)\,.
\]
It is convenient to use standard plane-wave solutions of the Klein-Gordon
equation
\[
f_\pp=e^{i\omega_\pp t-i\pp\cdot\xx}\,,\qquad f^*_\pp=e^{-i\omega_\pp t+i\pp\cdot\xx}\,,
\qquad\text{where}\qquad \omega_\pp=\sqrt{\pp^2+m^2}\,.
\]
Denoting
\[
\hat A_\pp=\hat A_{f_\pp}\,,\qquad \hat A^\dagger_\pp=\hat A_{f^*_\pp}\,,
\]
we find the commutation relations
\[
[\hat A_\pp,\hat A^\dagger_{\pp'}]=(2\pi)^3\,2\omega_\pp\,\delta^{(3)}(\pp-\pp')\,;\qquad
[\hat A_\pp,\hat A_{\pp'}]=0\,;\qquad
[\hat A^\dagger_\pp,\hat A^\dagger_{\pp'}]=0\,.
\]
These are the \emph{creation and annihilation operators}. The above operators
$\hat A_\pp$ and $\hat A^\dagger_\pp$ differ in normalization from $a_\pp$ and
$a^\dagger_\pp$ used in PS (see \textbf{Sect.~2.3})
\[
\hat A_\pp=\sqrt{2\omega_\pp}\,a_\pp\,;\qquad \hat A^\dagger_\pp=\sqrt{2\omega_\pp}\,a^\dagger_\pp\,.
\]
Advantage of $A,A^\dagger$ is their Lorentz invariance.

The Hamiltonian can be expressed through the creation and annihilation
operators as
\begin{equation}\label{eq:HAA}
\hat H=\int\frac{d^3\pp}{(2\pi)^3}\,\frac14\,\Bigl(\hat A^\dagger_\pp\,\hat A_\pp+\hat A_\pp\,\hat A^\dagger_\pp\Bigr)\,,
\end{equation}
which indeed is the Hamiltonian of a collection of harmonic oscillators, one
oscillator per each wave-vector $\pp$, the associated frequency being
$\omega_\pp$. It is straightforward to check that $\hat A^\dagger_\pp$ raises
(and $\hat A_\pp$ lowers) the energy by the amount $\omega_\pp$, namely
\[
[\hat H,\hat A^\dagger_\pp]=\omega_\pp\,\hat A^\dagger_\pp\,,\qquad
[\hat H,\hat A_\pp]=-\,\omega_\pp\,\hat A_\pp\,.
\]
Similarly, for the momentum
\[
\hat{\PPv}=-\int d^3\xx\,\hat\pi(\xx)\,\nabla\hat\varphi(\xx)
\]
we have
\[
\hat{\PPv}=\int\frac{d^3\pp}{(2\pi)^3}\,\frac{\pp}{4\omega_\pp}\,\Bigl(\hat A^\dagger_\pp\,\hat A_\pp+\hat A_\pp\,\hat A^\dagger_\pp\Bigr)\,.
\]
Hence
\[
[\hat{\PPv},\hat A^\dagger_\pp]=\pp\,\hat A^\dagger_\pp\,,\qquad
[\hat{\PPv},\hat A_\pp]=-\,\pp\,\hat A_\pp\,,
\]
i.e.\ $\hat A^\dagger_\pp$ adds $\pp$ to the total momentum.

To find a physically acceptable solution of quantum KG theory it remains to
find appropriate representation of the above commutation relations in a
Hilbert space $\Ham$. Classically, the Hamiltonian \eqref{eq:HKG} is
non-negative definite. Correspondingly, under physically meaningful
quantization the quantum Hamiltonian \eqref{eq:HAA} must be bounded from
below. Recall from the theory of harmonic oscillator that in order to satisfy
this property one has to assume that $\Ham$ contains the ground state
$\ket{0}$ such that
\[
\hat A_\pp\ket{0}=0\qquad\text{for all}\quad\pp\,.
\]
Just as in the harmonic oscillator, by applying the operators
$A^\dagger_\pp$ to $\ket{0}$ one generates the complete space of normalizable
states. In QFT this is referred as the \emph{Fock space},
\[
\Ham_{\mathrm{KG}}=\text{Fock space}=\operatorname{Span}\Bigl\{\hat A^\dagger_{\pp_1}\hat A^\dagger_{\pp_2}\cdots\hat A^\dagger_{\pp_N}\ket{0}\Bigr\}\,.
\]
The state $\ket{0}$ is called \emph{vacuum state}, or \emph{Fock vacuum}. Note
that the above basis already solves the diagonalization problem for $\hat H$:
these vectors are the eigenstates of both $\hat H$ and $\hat{\PPv}$, with the
eigenvalues
\begin{align*}
E&=\omega_{\pp_1}+\omega_{\pp_2}+\ldots+\omega_{\pp_N}+E_0\,,\\
\PPv&=\pp_1+\pp_2+\ldots+\pp_N+\PPv_0\,,
\end{align*}
where $E_0$ and $\PPv_0$ are the eigenvalues associated with the vacuum state
$\ket{0}$,
\[
\hat H\ket{0}=E_0\ket{0}\,,\qquad \hat{\PPv}\ket{0}=\PPv_0\ket{0}\,.
\]
The basic vectors $\hat A^\dagger_{\pp_1}\hat A^\dagger_{\pp_2}\cdots\hat A^\dagger_{\pp_N}\ket{0}$
are interpreted as the states of $N$ identical free (Bose) particles, with
the momenta $\pp_1,\dots,\pp_N$.

The above expression for the Hamiltonian can be rewritten as
\[
\hat H=E_0+\int\frac{d^3\pp}{(2\pi)^3}\,\frac12\,\hat A^\dagger_\pp\,\hat A_\pp\,,
\qquad\text{with}\qquad
E_0=\int\frac{d^3\pp}{(2\pi)^3}\,\frac14\,[\hat A_\pp,\hat A^\dagger_\pp]\,,
\]
which of course is the sum of the zero-point energies of all the constituent
oscillators.

There are two evident problems with the last expression:

(i) It is not zero, and moreover it contains seemingly meaningless factor
$\delta^{(3)}(\mathbf{0})$ (recall that
$[\hat A_\pp,\hat A^\dagger_{\pp'}]=2\omega_\pp\,(2\pi)^3\,\delta^{(3)}(\pp-\pp')$).

(ii) The remaining integral over $\pp$ diverges at large $\pp$.

Note that these problems don't seem to be as bad for $\PPv_0$, where we have
\[
\PPv_0=\int\frac{d^3\pp}{(2\pi)^3}\,\frac{\pp}{4\omega_\pp}\,[\hat A_\pp,\hat A^\dagger_\pp]=0\,,
\]
since the contributions of $\pp$ and $-\pp$ cancel pairwise. This is good, as
we do not want the vacuum to have any net momentum. But this seems to create
another problem, if one thinks that the energy and momentum should be the
components of 4-vector $(E_0,\PPv_0)$. If $\PPv_0=0$ and $E_0\neq0$, it seems
the Lorentz invariance is broken.

To sort this part of the problem out, let us recall that we are dealing with
a system in which the degrees of freedom $\varphi(\xx)$ are attached to all
points of the space. For such system one expects that its vacuum energy
$E_0$, whatever it is, must be proportional to the volume of the space,
\[
E_0=\varepsilon\,V^{(3)}\,,
\]
where $\varepsilon$ is the vacuum energy density. This answers the question
how to interpret the $\delta^{(3)}(\mathbf 0)$ factor: Recall that
\[
(2\pi)^3\,\delta^{(3)}(\pp-\pp')=\int d^3\xx\,e^{i(\pp-\pp')\cdot\xx}\,,
\]
and so the above nonsense factor in fact is
\[
(2\pi)^3\,\delta^{(3)}(\mathbf 0)=\int d^3\xx\ \to\ V^{(3)}\,.
\]
This also solves the problem of Lorentz invariance. The energy density
$\varepsilon$ transforms as $00$ component of 4-tensor (the energy-momentum
tensor),
\[
\varepsilon=\bra{0}T^{00}\ket{0}\,.
\]
The Lorentz-invariance of the vacuum state $\ket{0}$ only requires that
\[
\bra{0}T^{\mu\nu}\ket{0}=-\varepsilon_0\,\eta^{\mu\nu}\,,\qquad
\eta^{\mu\nu}=\operatorname{diag}(-1,+1,+1,+1)\,,
\]
and we soon will have techniques to check that this equation indeed holds in
the KG theory. It is OK for a Lorentz-invariant state to have nonzero energy
density.

Still, something has to be done about the badly divergent integral for this
energy density,
\[
\varepsilon_0=\frac12\int\frac{d^3\pp}{(2\pi)^3}\,\omega_\pp\,.
\]
The integral diverges at large momenta, i.e.\ the divergent contribution comes
from small length scales. This is not too surprising, since our idea of the
field theory was to associate a degree of freedom with every point of the
space, and there are continuously many points in each finite part of it. One
could assume that perhaps some microscopic physics takes care of the problem,
replacing the above mathematically meaningless expression by the ``cutoff''
integral
\[
\varepsilon_0=\frac12\int\frac{d^3\pp}{(2\pi)^3}\,\omega_\pp\,\Phi\left(\frac{\pp^2}{\Lambda^2}\right)\,,
\]
where the ``cutoff'' factor $\Phi(\xi)$ is close to 1 when $\xi\ll1$ and
decays fast at large $\xi$, and $\Lambda$ represents the energy scale beyond
which the unknown microscopic physics becomes essential.

Whereas introducing the cutoff rids us of mathematically meaningless
expression, needless to say it does not solve the problem of vacuum energy
density, which now is finite but completely undetermined. In many practical
cases one can simply ignore the problem. Large class of physical quantities
are only sensitive to differences of energies, and in such problems the
vacuum energy is invisible. But in fact the vacuum energy is measurable, in
principle. Generally, if we want to be able to apply QFT to physics, we need
to know what happens when we couple it to gravity, and gravity is very
sensitive to the energy density. Even without gravity, suppose we can confine
our system to finite part of the space, say inside a box. Then we want to know
what forces it exerts upon the wall of that box. The last situation can
actually be realized in laboratory, and it leads to a measurable effect, the
so called Casimir effect. We are going to discuss it next, but let me make one
more remark, in order to clarify the nature of the problem.

When a divergent integral appears in otherwise sensible calculations, it
usually signals some ambiguity in the theory. Indeed, there is an ambiguity
in the definition of the KG theory. One can add a constant term to the
Lagrangian density,
\[
\Lag\to\Lag+\Const\,.
\]
This extra term is fully consistent with both relativistic symmetry and
locality. Of course this modification is irrelevant in the classical theory
(again, as long as gravity is ignored), since it does not affect the
equations of motion. But the Lagrangian density will pass this term to the
energy-momentum tensor, and finally to the Hamiltonian,
\[
\hat H\to\hat H-\Const\int d^3\xx\,.
\]
One can now ``absorb'' the vacuum energy divergence into this constant. We
write
\[
\Const=\varepsilon_0-\varepsilon\,,
\]
where $\varepsilon_0$ is the above ugly integral, and $\varepsilon$ is a
finite constant (of the dimension $[\text{Mass}]^4$). The result is that the
divergent integral disappears, but at the same time we learn that the quantum
KG theory in fact has not one but two parameters, $m^2$ and $\varepsilon$.
This extra term we added to the action is the simplest example of what we
will call the \emph{counterterms}.
